% This work is made available under the terms of the
% Creative Commons Attribution-ShareAlike 4.0 license,
% http://creativecommons.org/licenses/by-sa/4.0/.

\documentclass[a4paper]{book}

\usepackage{wrapfig}
\usepackage{graphicx}
\usepackage{hyperref}
\usepackage{multirow}
\usepackage{scalefnt}
\usepackage{tikz}

% watermark -- for draft stage
%\usepackage[firstpage]{draftwatermark}
%\SetWatermarkLightness{0.9}
%\SetWatermarkScale{5}

% Copyright (c) 2009 by the University of Waikato, Hamilton, NZ. 
% This work is made available under the terms of the 
% Creative Commons Attribution-ShareAlike 4.0 license,
% http://creativecommons.org/licenses/by-sa/4.0/.
%
% Version: $Revision: 5479 $

\newenvironment{tight_itemize}{
\begin{itemize}
  \setlength{\itemsep}{1pt}
  \setlength{\parskip}{0pt}
  \setlength{\parsep}{0pt}}{\end{itemize}
}

\newenvironment{tight_enumerate}{
\begin{enumerate}
  \setlength{\itemsep}{1pt}
  \setlength{\parskip}{0pt}
  \setlength{\parsep}{0pt}}{\end{enumerate}
}

% if you just need a simple heading
% Usage:
%   \heading{the text of the heading}
\newcommand{\heading}[1]{
  \vspace{0.3cm} \noindent \textbf{#1} \newline
}

\newcommand{\icon}[1]{\tikz[baseline=-3pt]\node[inner sep=0pt,outer sep=0pt]{\includegraphics[height=1.1em]{#1}};}


\title{
  \textbf{ADAMS} \\
  {\Large \textbf{A}dvanced \textbf{D}ata mining \textbf{A}nd \textbf{M}achine
  learning \textbf{S}ystem} \\
  {\Large Module: adams-rabbitmq} \\
  \vspace{1cm}
  \includegraphics[width=2cm]{images/rabbitmq-module.png} \\
}
\author{
  Peter Reutemann
}

\setcounter{secnumdepth}{3}
\setcounter{tocdepth}{3}

\begin{document}

\begin{titlepage}
\maketitle

\thispagestyle{empty}
\center
\begin{table}[b]
	\begin{tabular}{c l l}
		\parbox[c][2cm]{2cm}{\copyright 2019} &
		\parbox[c][2cm]{5cm}{\includegraphics[width=5cm]{images/coat_of_arms.pdf}} \\
	\end{tabular}
	\includegraphics[width=12cm]{images/cc.png} \\
\end{table}

\end{titlepage}

\tableofcontents
%\listoffigures
%\listoftables

%%%%%%%%%%%%%%%%%%%%%%%%%%%%%%%%%%%
\chapter{Introduction}
RabbitMQ\cite{rabbitmq} is an open source message broker. It supports:
\begin{tight_itemize}
  \item multiple messaging protocols
  \item message queuing
  \item delivery acknowledgement
  \item flexible routing to queues
  \item multiple exchange type
\end{tight_itemize}
ADAMS makes use of it for publishing and consuming messages, as well as for
remote procedure calls (RPC).

%%%%%%%%%%%%%%%%%%%%%%%%%%%%%%%%%%%
\chapter{Installation}
Installation instructions for all available platforms can be found here: \\
\url{https://www.rabbitmq.com/download.html}{}

\section{Linux}
You can install the RabbitMQ server on Debian systems like this:
\begin{verbatim}
sudo apt-get install rabbitmq-server
\end{verbatim}

\noindent For RPM-based distributions, please follow the instructions here: \\
\url{https://www.rabbitmq.com/install-rpm.html}{}

\section{Windows}
For installing the server on Windows, please follow the instructions here: \\
\url{https://www.rabbitmq.com/install-windows.html}{}

\section{TLS Support}
\label{tls_support}
You can find information on how to set up RabbitMQ using TLS on the
following URL: \\
\url{https://www.rabbitmq.com/ssl.html}{}


%%%%%%%%%%%%%%%%%%%%%%%%%%%%%%%%%%%
\chapter{Flow}
The following standalones are available:
\begin{tight_itemize}
  \item \textit{DataExchangeServerConnection} -- defines a data exchange server connection.
  \item \textit{RabbitMQConnection} -- defines a RabbitMQ broker connection.
  \item \textit{RabbitMQChannelAction} -- peforms the specified channel action.
\end{tight_itemize}

\noindent The following sources are available:
\begin{tight_itemize}
  \item \textit{RabbitMQConsume} -- outputs the data it consumes from either
  a queue or an exchange.
\end{tight_itemize}

\noindent The following transformers are available:
\begin{tight_itemize}
  \item \textit{DataExchangeServerDownload} -- downloads data from a data exchange
  server using the specified token.
  \item \textit{DataExchangeServerRemove} -- removes data from a data exchange
  using the received token.
  \item \textit{DataExchangeServerUpload} -- uploads a file or byte array to a data exchange
  server and forwards the received token.
  \item \textit{RabbitMQRemoteProcedureCall} -- for performing remote procedure
  calls via a RabbitMQ broker.
\end{tight_itemize}

\noindent The following sinks are available:
\begin{tight_itemize}
  \item \textit{RabbitMQMessageDeliveryAction} -- for executing an action using
  the incoming delivery tag, e.g., acknowledging it.
  \item \textit{RabbitMQPublish} -- publishes the incoming data to a queue
  or an exchange.
\end{tight_itemize}

\noindent The following control actors are available:
\begin{tight_itemize}
  \item \textit{RabbitMQRemoteSubProcess} -- executes the sub-actors remotely
  for processing the input, forwards the generated output.
  \item \textit{RabbitMQRemoteTee} -- executes the sub-actors remotely
  for processing the input, forwards the input.
  \item \textit{RabbitMQRemoteTrigger} -- executes the sub-actors remotely,
  forwards the input.
\end{tight_itemize}

\noindent The following conversions are available:
\begin{tight_itemize}
  \item \textit{RabbitMQEnvelopeToMap} -- converts RabbitMQ envelope into a
  Java Map, e.g., to obtain a delivery tag to acknowledge consumption of a message.
  \item \textit{RabbitMQPropertiesToMap} -- converts RabbitMQ basic properties
  into a Java Map, e.g., to obtain the reply-to queue name.
\end{tight_itemize}

\noindent The following JobRunners are available:
\begin{tight_itemize}
  \item \textit{RabbitMQJobRunner} -- distributes the jobs using the specified
  broker.
\end{tight_itemize}

\section{TLS Support}
If you want to take advantage of TLS, i.e., encrypting your connections to
the RabbitMQ server, then you need to configure your server for that (see
\ref{tls_support}) and encapsulate your actual connection factory used by
\textit{RabbitMQConnection} in one of the following meta-factories:
\begin{tight_itemize}
  \item \textit{NonValidatingSSLConnectionFactory} -- only to be used for testing/development
  \item \textit{SSLConnectionFactory}
\end{tight_itemize}

\section{Large payloads}
In theory, you can use RabbitMQ to send any amount of data. However, queues
may be constrained in disk size and sending large amounts (e.g., processing
large datasets) will also slow down the overall performance of the message broker.

As long as the sending and receiving ends have access to the same shared directory,
it is possible to exchange the actual payload via this directory and only send
the filename via the message broker. For this scenario, you can use the following
converters:
\begin{tight_itemize}
  \item adams.core.net.rabbitmq.send.FileBasedConverter
  \item adams.core.net.rabbitmq.receive.FileBasedConverter
\end{tight_itemize}
The receiving converter automatically deletes the payload file ones read
successfully.

In case a shared directory is not an option, see chapter \ref{dex} on how to use
a \textit{data exchange server} instead.


%%%%%%%%%%%%%%%%%%%%%%%%%%%%%%%%%%%
\chapter{Data exchange server}
\label{dex}
The \textit{Data exchange server} (DEX) is a simple REST webservice plugin that enables
one to store data temporarily on a server and access it. The underlying idea is
to avoid clogging up message queues on RabbitMQ servers. Rather than sending the
data as well in the message, only a token that is associated with the actual
data on such an exchange server will get sent through RabbitMQ.

\noindent The following \textit{RESTProvider} that can be used for this service:
\begin{verbatim}
adams.flow.rest.dex.DataExchangeServer
\end{verbatim}

\noindent The \textit{RESTPlugin} itself (used internally by the above provider)
is the following:
\begin{verbatim}
adams.flow.rest.dex.DataExchange
\end{verbatim}

\noindent You can use a data exchange server for offloading payloads from
RabbitMQ messages with the following converters:
\begin{tight_itemize}
  \item \textit{adams.core.net.rabbitmq.receive.DataExchangeServerBasedConverter}
  \item \textit{adams.core.net.rabbitmq.send.DataExchangeServerBasedConverter}
\end{tight_itemize}

\noindent The following example flows demonstrate the data exchange functionality:
\begin{tight_itemize}
  \item \textit{adams-rabbitmq-dex\_server.flow} -- starts an in-memory data exchange server.
  \item \textit{adams-rabbitmq-dex\_upload.flow} -- uploads a file to the data exchange server.
  \item \textit{adams-rabbitmq-dex\_download.flow} -- downloads data from the data exchange server using the entered token.
  \item \textit{adams-rabbitmq-dex\_queue\_publish.flow} -- publishes data through RabbitMQ via a data exchange server.
  \item \textit{adams-rabbitmq-dex\_queue\_consume.flow} -- consumes data through RabbitMQ via a data exchange server.
\end{tight_itemize}

%%%%%%%%%%%%%%%%%%%%%%%%%%%%%%%%%%%
% Copyright (c) 2009-2012 by the University of Waikato, Hamilton, NZ. 
% This work is made available under the terms of the 
% Creative Commons Attribution-ShareAlike 4.0 license,
% http://creativecommons.org/licenses/by-sa/4.0/.
%
% Version: $Revision$

\begin{thebibliography}{999}
	% to make the bibliography appear in the TOC
	\addcontentsline{toc}{chapter}{Bibliography}

    % references
	\bibitem{adams}
		\textit{ADAMS} -- Advanced Data mining and Machine learning System \\
		\url{https://adams.cms.waikato.ac.nz/}{}
		
	\bibitem{heatmap}
		\textit{Heat map} -- WikiPedia article \\
		\url{http://en.wikipedia.org/wiki/Heat_map}{}

\end{thebibliography}


\end{document}
